\section{Recommendations and Outlook}
\label{sec:recommendations}

The results above are not an argument against using earnings calls. They are an
argument about what an evaluation has to control for before a content effect can
be claimed. We suggest that future work on this task report, at minimum:
(1)~a \textbf{company-disjoint} split next to the temporal one, since identity
carries most of the apparent signal; (2)~a \textbf{gap} of at least the longest
horizon between training and test dates; (3)~\textbf{past-volatility baselines}
(persistence, HAR-RV) and, on change targets, a constant, as the reference for
every claim; (4)~the \textbf{same-company and random swaps} or the released
audit command, so that a gain can be attributed to the call rather than to the
company; (5)~\textbf{day~0 measured} from the release time rather than the call
date; and (6)~the \textbf{horizon definition} (calendar or trading days) stated
with every number. \ecvol-bench provides each of these for FinCall-Surprise and
MAEC, and the audit command applies the checks to any prediction file.

There is also reason to expect that content signal can be found once it is
looked for in the right place. Measured release times make
\textbf{event-window targets} possible, such as the volatility of the first
sessions after the release, where language about guidance and surprise should
matter most and where same-company history explains least. \textbf{Structured
features} extracted from the call, such as guidance changes and the gap between
the quantitative surprise and management's framing of it, can be audited one
field at a time, in contrast to pooled embeddings, and can be combined with the
fast numeric signal that \citet{fastnumbers2026} show arrives first.
\textbf{Audio} has so far been tested only on lossy webcast transcodes without
speaker diarisation; recordings at source quality with resolved speakers are the
fair test of prosody. \textbf{Larger company-disjoint corpora} are now within
reach, since MM-FinEval \citep{mmfineval2026} ships ticker and date fields and the
checks here transfer to it directly. Each of these can be evaluated against the
released baselines on the first day, which is what the benchmark is for.
